% Reconstructed LaTeX source for the four appended pages of
% papers/Unified_Intelligence_Theory_and_AI_Level_v2_6.pdf (pages 66-69).
%
% Those pages were NOT produced by LaTeX. They were generated outside the TeX
% source and merged onto the 65-page pdfTeX build with pypdf; pages 1-65 use
% Nimbus Roman (pdfTeX), pages 66-69 use Noto Sans. No editable source for them
% was found in the repositories or release bundles that were searched.
%
% This file is a faithful transcription of the published text of those four
% pages into LaTeX, so that the v2.6 appendix survives the move to a TMLR
% template. It is a reconstruction, not the original source, and the wording is
% the PDF's own.
%
% Both addenda have since been folded into the main manuscript:
% main.tex carries them as the section "Cumulative structure is typed, not
% universal" (Table tab:typed) and the surrounding text, so a submission built
% from main.tex should normally NOT append this file. Keep it for provenance,
% or \input it if the appendix form is wanted verbatim.
%
% Build on its own with: pdflatex addenda_v2_5_v2_6.tex

\documentclass[11pt]{article}
\usepackage[letterpaper,margin=1in]{geometry}
\usepackage[T1]{fontenc}
\usepackage[utf8]{inputenc}
\usepackage{booktabs}
\usepackage{tabularx}
\usepackage{amsmath,amssymb}
\usepackage{parskip}
\sloppy
\renewcommand{\arraystretch}{1.3}

\newcommand{\addendumtitle}[2]{%
  \begin{center}
    {\large\bfseries Unified Intelligence Theory and the Artificial Intelligence Level}\\[0.4em]
    {\bfseries #1}\\[0.2em]
    {\small #2}
  \end{center}
  \vspace{0.5em}}

\begin{document}

%% ---------------------------------------------------------------- v2.5 ----
\addendumtitle{v2.5 Revision Addendum}{AI-Level Naming, Uniform Testing Grammar, and Dataset Binding}

\section*{1. Authority and current naming}

This revision treats Unified Intelligence Theory as the semantic basis of the measurement system. The current framework name is AI-Level; \texttt{AI\_Level} is the file-system form. New papers, dataset packages, schemas, commands, protocol identifiers and result records use AI-Level naming. The change of name does not create a new intelligence family and does not change the meanings of $C$, $I$, $O$, $DI$, $SA$, $M$, $T$, $H$, $HG$ or $U$.

The ontology remains:

\begin{quote}
Conceptual core $= [C, I, O, DI, SA]$; measured core $= [C, I, O, T, H, SA]$, with $DI$ derived from the delivered-outcome surface over $(T, H, p)$. Memory Capability $M$ is independently certified and joined to $I$ only through one-way prerequisites. GUI/screen cognition is a domain of $C$ with GP0--GP5 as a diagnostic perceptual subscale. $HG$ is an engineering/conformance ladder. $U$ is the cumulative Unified gate.
\end{quote}

\section*{2. One experimental grammar for every ladder}

Every level, band and rung is now specified using AI-Level-Method-1.0, the same causal and retention-oriented style used for Individual Intelligence. This is a measurement grammar, not a claim that all objects are the same kind of intelligence.

\small
\begin{tabularx}{\textwidth}{@{}>{\raggedright\arraybackslash}p{0.17\textwidth}>{\raggedright\arraybackslash}X@{}}
\toprule
Field & Requirement \\
\midrule
1 Construct & State the capability as an invariant contrast or structural property. \\
2 Prerequisite & Name what can block the claim but can never promote it by itself. \\
3 Intervention & Specify the manipulation, lifecycle discontinuity, task structure, or role/resource perturbation. \\
4 Ground truth & Identify hidden or externally grounded state outside the candidate write set. \\
5 Witness + control & Run a matched control/ablation whenever causal attribution is possible. \\
6 Metrics & Record primary outcomes and separated failure modes. \\
7 Factor gate & Certification is a product/conjunction of named auditable factors. \\
8 Retention & Require all lower levels relevant to the ladder. \\
9 Reliability & Apply preregistered repeated-sampling/uncertainty rules. \\
10 Binding & Mark bound, development-bound, or specification-only. \\
\bottomrule
\end{tabularx}
\normalsize

\subsection*{Generic certification form}

\begin{quote}
$A \models d_k$ iff $V_{\mathrm{new},d,k}(A) = 1$ \textsc{and} $K_{d,<k}(A) = 1$, with any independent prerequisite gate included conjunctively. Continuous achievement may summarize progress but never overrides a failed ordinal gate.
\end{quote}

\section*{3. Ladder-specific interpretation}

\small
\begin{tabularx}{\textwidth}{@{}>{\raggedright\arraybackslash}p{0.07\textwidth}>{\raggedright\arraybackslash}X>{\raggedright\arraybackslash}X@{}}
\toprule
Object & What is tested in the common grammar & Key control \\
\midrule
$C$ & Depth of problem solving: reactive $\to$ contextual $\to$ compositional $\to$ strategic $\to$ expert $\to$ discovery $\to$ cross-domain $\to$ reach expansion. & Wrong-method, distractor, hypothesis uniqueness, domain ablation, or instrument ablation. \\
$I$ & Persistence, learning, governed $\Theta$ change, $\Psi$ improvement, persistent discovery, reach expansion. & Restart/no-state, experience/ablated, pre/post mechanism, meta-campaign, incorporation control, instrument ablation. \\
$M$ & Working state $\to$ persistence $\to$ provenance $\to$ consolidation $\to$ management $\to$ lineage $\to$ memory-mechanism evolution. & State ablation, raw-episode control, seeded corruption, $\Phi_0/\Phi_1$ paired evaluation. \\
$O$ & Routing $\to$ persistent organization $\to$ adaptive routing $\to$ workflow self-improvement $\to$ recursive organization $\to$ collective discovery $\to$ open-ended collective reach. & Role/state ablation, fixed-workflow control, same-members pre/post organization, organizer-process control. \\
$SA$ & No grounded model $\to$ state $\to$ limits $\to$ causal diagnosis $\to$ self-change prediction $\to$ self-experiment $\to$ mission/role $\to$ reflexive evolving self-model. & Stale-state decoy, permission twin, fault intervention, prediction/post-change comparison, experiment execution, role perturbation, lineage audit. \\
$T$ & Structural property of the task from atomic through open-ended charter; difficulty is in the task, not inferred from who solves it. & Admissibility test showing required dependency/uncertainty/mission structure is present. \\
$H$ & Externally classified human cognitive intervention from H0 to H5 using the intervention ledger. & Independent locus classification; same task under differing human assistance where feasible. \\
$DI$ & Delivered-outcome frontier over $(T, H, p)$. & Held-back/refusal arm must differ from false completion; lower-$T$ retention. \\
$HG$ & Harness mechanism exists, fires, and causally changes the same-seed outcome compared with prior rung. & Same model, same seeds, previous-rung matched control. \\
$U$ & Conjunction of the required coordinate gates plus lower-$U$ retention. & No averaging can compensate a failed bottleneck. \\
\bottomrule
\end{tabularx}
\normalsize

\section*{4. Operational Self-Awareness is fully grounded}

$SA$ is defined operationally. Fluent self-description is never sufficient. Each rung must be scored against external runtime truth or an intervention whose effect the tested system does not control.

\small
\begin{tabularx}{\textwidth}{@{}>{\raggedright\arraybackslash}p{0.17\textwidth}>{\raggedright\arraybackslash}X@{}}
\toprule
Level & Canonical evidence \\
\midrule
SA0 & Baseline category: no higher grounded self-model is certified. \\
SA1 & Randomize identity/role/phase/resources; query grounded state with stale-note decoys. \\
SA2 & Randomize tools, permissions, authority, freshness and unavailable information; require correct capability/limit behavior. \\
SA3 & Inject one of several controlled faults; require mechanism-level diagnosis and validate by repair/ablation/counterfactual. \\
SA4 & Before a governed self-change, preregister predicted gains/regressions; compare with frozen post-change results. \\
SA5 & Present ambiguous self-unknowns; require a discriminating self-experiment, execute it, and score the belief update. \\
SA6 & Across projects and role/authority changes, maintain current mission, responsibility and relationship state; suppress stale roles. \\
SA$\Omega$ & Across validated cognitive/objective evolution, maintain a current provenance-linked self-model and suppress superseded self-beliefs. \\
\bottomrule
\end{tabularx}
\normalsize

$SA$ is cumulative for certification: SA$\Omega \to$ SA6 $\to$ SA5 $\to$ SA4 $\to$ SA3 $\to$ SA2 $\to$ SA1. SA0 is the baseline rather than a positive promotion claim.

\section*{5. GUI cognition remains inside Cognitive Intelligence}

Method B is retained: $C^{\mathrm{GUI}}_k$ is a GUI-grounded witness of the unchanged $C_k$ semantic level, while GP0--GP5 measures perceptual screen representation. GP does not map one-to-one to $C$ and cannot promote $C$ by itself. Native-screen, oracle-screen and perfect-actuator arms separate perception, cognition and action failures.

\section*{6. Dataset binding}

Every currently defined ladder now has a machine-readable canonical method entry and public development witness in AI-Level Benchmark Dataset v0.2. Upper open-ended and expensive longitudinal forms may remain specification-only until secure runners exist; they are still represented explicitly and never imputed as zero.

\subsection*{Dataset identifiers}

\begin{quote}
Dataset: AI-Level Benchmark Dataset v0.2; protocol: AI-Level-Method-1.0; benchmark basis: AI-Level Bench v0.2; theory basis: Unified Intelligence Theory and the Artificial Intelligence Level v2.5.
\end{quote}

\clearpage

%% ---------------------------------------------------------------- v2.6 ----
\addendumtitle{v2.6 Cumulative-Structure Clarification}{Capability Ladders, GP, Delegation Frontier, and Harness Rungs}

\section*{1. Cumulative structure is typed, not universal}

The framework retains one experimental grammar for all measured objects, but it does not treat every ordered object as the same kind of ladder. This revision distinguishes cumulative capability, cumulative diagnostic capability, cumulative frontier, ordered measurement axes, and cumulative engineering structure. The purpose is to preserve the cumulative semantics of intelligence while avoiding false inheritance for task difficulty $T$ and human intervention $H$.

\small
\begin{tabularx}{\textwidth}{@{}>{\raggedright\arraybackslash}p{0.10\textwidth}>{\raggedright\arraybackslash}p{0.17\textwidth}>{\raggedright\arraybackslash}X>{\raggedright\arraybackslash}X@{}}
\toprule
Object & Type & Cumulative rule & Why \\
\midrule
$C$ & Capability ladder & Hard cumulative & A higher cognitive level must retain lower cognitive capabilities. \\
$I$ & Capability ladder & Hard cumulative & Persistent learning/self-improvement claims require lower Individual capabilities. \\
$M$ & Supporting capability ladder & Hard cumulative & Higher memory management cannot bypass persistence/provenance/consolidation. \\
$O$ & Capability ladder & Hard cumulative & Higher collective improvement requires retained coordination and organizational state. \\
SA1--SA$\Omega$ & Capability ladder & Hard cumulative & Deeper self-modeling must retain grounded state/capability/causal awareness. \\
SA0 & Baseline category & Not retained & SA0 denotes absence of a reliable grounded self-model; SA1+ cannot literally retain that absence. \\
GP0--GP5 & Diagnostic capability sub-ladder & Hard cumulative within GP & Objects $\to$ grounding $\to$ state $\to$ dynamics $\to$ novel-interface generalization build on earlier perceptual abilities. \\
$DI$ & Delegation surface/frontier & Cumulative along $T$ at fixed $H, p$ & A hard-band pass cannot leap over a failed easier band. \\
$T$ & Task-classification axis & Not cumulative by itself & $T$ describes structural properties of tasks, not abilities of the system. \\
$H$ & Human-intervention axis & Not cumulative by itself & $H$ classifies assistance supplied; lower $H$ is less intervention. \\
$HG$ & Engineering reference ladder & Structurally cumulative & Each reference rung is a strict mechanism superset so causal harness gain is attributable. \\
$U$ & Gated integration ladder & Hard cumulative & Every Unified promotion retains all lower $U$ gates and all required coordinate certificates. \\
\bottomrule
\end{tabularx}
\normalsize

\section*{2. Revised general cumulative law}

For a cumulative capability ladder $X \in \{C, I, M, O, SA_{1:\Omega}, GP, U\}$, the level-$k$ certificate requires the new capability plus retention of every required lower capability:

\begin{quote}
$K_{X,<k}(A) = \prod_{j<k} K_{X,j}(A)$; $A \models X_k$ iff $V_{\mathrm{new},X,k}(A) = 1$ \textsc{and} $K_{X,<k}(A) = 1$, with any independent cross-ladder prerequisite added conjunctively.

For continuous achievement, $q^{*}_{X,k} = \min_{j \le k} q_{X,j}$. This prevents a high-level witness from compensating for a failed lower-level retention requirement.
\end{quote}

The global wording is therefore sharpened to: \emph{No certified capability level may skip required lower capabilities.} The sentence is not applied literally to $T$ or $H$, because they are not capability ladders.

\section*{3. Cognitive Intelligence $C$ remains cumulative}

$C0$--$C\Omega$ is a hard cumulative capability ladder. A system that passes a C5 discovery witness but fails C2 composition has evidence of a C5-type behavior, not a C5 certificate. Domain-specific cognition, including $C^{\mathrm{GUI}}$, inherits the ordinary $C$ retention gate rather than creating a second semantic $C$ ladder.

\begin{quote}
For cognitive domain $d$: $q^{*}_{C,k,d} = \min_{j \le k} q_{C,j,d}$. Domain aggregation is performed after within-domain level retention, under the declared domain policy and floors.
\end{quote}

\section*{4. GP becomes a cumulative diagnostic capability sub-ladder}

GUI Perception (GP) remains inside the GUI/screen domain of $C$ and remains diagnostic: GP evidence reports the fidelity and generality of the screen representation supplied to cognition. The semantic containment remains $GP \subset C^{\mathrm{GUI}} \subset C$, with $GP_g \not\equiv C_g$. This revision adds only the cumulative rule inside GP.

\begin{quote}
$GP5 \Rightarrow GP4 \Rightarrow GP3 \Rightarrow GP2 \Rightarrow GP1 \Rightarrow GP0$

$A \models GP_g$ iff $V_{GP,g}(A) = 1$ \textsc{and} $K_{GP,<g}(A) = 1$; $q^{*}_{GP,g} = \min_{j \le g} q_{GP,j}$.
\end{quote}

\small
\begin{tabularx}{\textwidth}{@{}>{\raggedright\arraybackslash}p{0.07\textwidth}>{\raggedright\arraybackslash}X>{\raggedright\arraybackslash}p{0.27\textwidth}@{}}
\toprule
GP & New capability at this rung & Required retained GP capabilities \\
\midrule
GP0 & Pixel/text/basic visual recognition & none \\
GP1 & UI-element recognition & GP0 \\
GP2 & Spatial grounding and hierarchy & GP0--GP1 \\
GP3 & Interface-state understanding & GP0--GP2 \\
GP4 & Dynamic GUI understanding across frames/actions & GP0--GP3 \\
GP5 & Novel-interface perceptual generalization & GP0--GP4 \\
\bottomrule
\end{tabularx}
\normalsize

GP retention is internal to GP. It does not produce the implications $GP_g \Rightarrow C^{\mathrm{GUI}}_g$, $GP_g \Rightarrow C_g$, or $GP_g \Rightarrow U_g$. A native-image $C^{\mathrm{GUI}}$ witness may declare a minimum GP prerequisite $\mu(k)$; that prerequisite can block the native witness but can never promote $C$. The exact $\mu(k)$ remains form-specific and must be predeclared/calibrated.

\section*{5. Delegation is cumulative as a frontier, not as a $T$ ladder}

$T0$--$T\Omega$ are task bands defined by structural task properties. They do not inherit one another as task objects. $H0$--$H5$ are intervention classes and likewise do not carry lower-level retention. Cumulative behavior enters only when the system delegation frontier is certified at fixed intervention budget $h$ and reliability $p$:

\begin{quote}
$DF(h,p) = \max \{\, T_k : \min_{j \le k} S_{\mathrm{net}}(A, T_j, H \le h) \ge p \,\}$.
\end{quote}

Thus a successful T4 episode cannot certify a T4 delegation frontier if T2 or T3 fails at the same intervention ceiling and reliability. $T$ itself remains an ordered classification axis; $H$ remains an ordered assistance axis.

\section*{6. $HG$ is cumulative engineering structure; $U$ is cumulative intelligence integration}

The standardized reference harness ladder uses strict structural nesting: $HG_g = HG_{g-1} \cup \Delta_g$. The previous rung on the same model and seeds is the causal control. This is cumulative architecture, not intelligence: an HG4-conformant harness can still instantiate an I2 system.

$U$ remains hard cumulative. $U_n$ is issued only when all required coordinate certificates, the delegation gate, and $K_{U,<n}$ pass. No continuous score can compensate a failed $U$ bottleneck.

\section*{7. Consequence for AI-Level measurement and datasets}

Canonical records must declare \texttt{cumulative\_type}, \texttt{required\_lower\_levels}, \texttt{retention\_rule}, and no-promotion boundaries. GP forms are now represented as coordinate GP with \texttt{parent\_coordinate} $C$ and \texttt{parent\_domain} GUI/screen. $T$ and $H$ records explicitly declare non-cumulative axis semantics. $DI$ records alone apply cumulative $T$ retention at fixed $H, p$. This clarification changes the measurement contract, not the meanings of the underlying $C$, $I$, $M$, $O$, $SA$, $T$, $H$, $HG$, or $U$ labels.

\end{document}
